Die Fläche der Meere bleibt (fast) gleich, das Wasser kann sich nur nach oben ausdehnen. Für eine Wassersäule mit Querschnitt $A$ gilt $\Delta V = A\,\Delta h$ und $V_0 = A\,h$, also
\begin{align}
 \Delta h &= \frac{\Delta V}{A} = \frac{\gamma\,A\,h\,\Delta\vartheta}{A} = \dhMF
\end{align}
Mit $\gamma = \gaWaO$ (Wasser, FoTaBe S.~191):
\begin{align}
 \Delta h &= \gaWa\cdot\hMe\cdot\dTw = \dhM \approx \result{\dhMP{-2}{2}}
\end{align}
Das ist eine obere Abschätzung: Der Tabellenwert gilt für Wasser von \TtbO. Das Tiefenwasser der Ozeane ist nur wenige Grad warm, dort dehnt sich Wasser viel weniger aus. Zudem erwärmt sich vor allem die oberste Schicht.
